\chapter{سیستم‌های نوع محض لایه‌ای و ترجمه حذف وابستگی}\label{S2}
\pagestyle{plain}

\section{انگیزه و ایده اصلی}

\subsection{انگیزه}

در فصل پیشین، مفاهیم پایه‌ای حساب لامبدا، سیستم‌های کاهش انتزاعی، حساب لامبدا نوع‌دار ساده و سیستم‌های نوع محض را بررسی کردیم.
این چارچوب‌ها زمینه لازم برای درک مسائل پیشرفته‌تر نظریه انواع، به‌ویژه رابطه میان نرمال‌سازی ضعیف و قوی را فراهم آوردند.

سیستم‌های نوع محض، به‌عنوان چارچوبی یکپارچه برای توصیف سیستم‌های نوع مبتنی بر حساب لامبدا، بستری مناسب برای مطالعه خواص فرانظری نظیر نرمال‌سازی، یکتایی نوع و سازگاری منطقی فراهم می‌کنند.
با این حال، بررسی رابطه میان نرمال‌سازی ضعیف و نرمال‌سازی قوی در این سیستم‌ها، که در قالب حدس
\ref{BGK}
بیان می‌شود، در حضور انواع وابسته با دشواری‌های اساسی مواجه است.

مسئله اصلی از آنجا ناشی می‌شود که در سیستم‌های نوع وابسته، ساختار انواع می‌تواند به‌طور دلخواه به مقادیر عبارات وابسته باشد.
این وابستگی باعث می‌شود که رفتار مسیرهای کاهش
$\beta$
به‌شدت پیچیده شده و تحلیل آن‌ها، به‌ویژه در اثبات نرمال‌سازی قوی، با موانع جدی روبه‌رو گردد. در چنین شرایطی، روش‌های کلاسیک که برای سیستم‌های غیروابسته به‌کار می‌رود، مانند تکنیک‌های مبتنی بر کاهش‌پذیری یا استدلال‌های ترتیبی، دیگر به‌سادگی قابل اعمال نیستند.

\subsection{ایده اصلی}

\lr{Mull}
در رساله دکتری خود
\cite{mull2023}
بر این اصل استوار است که وابستگی‌های موجود در سیستم‌های نوع وابسته، لزوماً همگی برای بیان توان محاسباتی سیستم ضروری نیستند.
به بیان دیگر، بخشی از این وابستگی‌ها را می‌توان به صورت کنترل‌شده حذف یا محدود کرد، بی‌آنکه به خوش‌نوعی یا قابلیت بیان سیستم آسیبی وارد شود.

بر همین اساس،
\lr{Mull}
خانواده جدیدی از سیستم‌های نوع را تحت عنوان «سیستم‌های نوع محض لایه‌ای»\footnote{\lr{Tiered Pure Type Systems} یا به اختصار \lr{TPTS}}
معرفی می‌کند. در این سیستم‌ها، وابستگی میان انواع و عبارات به وسیله مفهوم «لایه» به طور صریح ساختاربندی می‌شود. لایه‌ها امکان تفکیک سطوح مختلف وابستگی را فراهم می‌کنند و از بروز وابستگی‌های چرخه‌ای یا کنترل‌نشده جلوگیری می‌نمایند.

مزیت کلیدی این لایه‌بندی آن است که امکان تعریف یک ترجمه حذف وابستگی از سیستم‌های نوع محض لایه‌ای به یک سیستم نوع محض فراهم می‌شود؛ به‌گونه‌ای که تحلیل نرمال‌سازی قوی در سیستم مقصد ساده‌تر بوده و می‌توان از آن برای استنتاج نرمال‌سازی قوی در سیستم مبدأ بهره گرفت.

در ادامه این فصل، ابتدا سیستم‌های نوع محض لایه‌ای به‌طور رسمی معرفی می‌شوند، سپس ترجمه حذف وابستگی و خواص اساسی آن مورد بررسی قرار می‌گیرد. این مباحث، زمینه لازم برای صوری‌سازی کامل این ترجمه و اثبات‌های مرتبط با آن در اثبات‌یار
\lr{Coq}
را که در فصل‌های بعدی ارائه می‌شوند، فراهم می‌سازند.

\section{سیستم‌های نوع محض لایه‌ای}

\subsection{لایه‌ها و شهود آن‌ها}

ایدهٔ اصلی در سیستم‌های نوع محض لایه‌ای، افزودن ساختاری صریح برای کنترل وابستگی‌ های نوعی است. این ساختار از طریق مفهوم «لایه» معرفی می‌شود. لایه‌ها سطوحی انتزاعی هستند که به متغیرها و عبارات نسبت داده می‌شوند و میزان وابستگی مجاز میان آن‌ها را محدود می‌کنند.

به‌طور شهودی، لایه‌ها را می‌توان به‌عنوان ابزاری برای مرتب‌سازی وابستگی‌ها در نظر گرفت؛ به‌گونه‌ای که یک عبارت در یک لایهٔ بالاتر مجاز نیست به‌طور دلخواه به عبارات لایه‌های پایین‌تر وابسته باشد. این محدودیت مانع از شکل‌گیری وابستگی‌های چرخه‌ای و کنترل‌نشده در سطح انواع می‌شود، که یکی از موانع اصلی در تحلیل نرمال‌سازی قوی در سیستم‌های وابسته است.

لازم به تأکید است که لایه‌ها با «رده‌ها» متفاوت‌اند. رده‌ها نقش طبقه‌بندی نحوی انواع و عبارات را ایفا می‌کنند، در حالی که لایه‌ها اطلاعاتی فرانظری درباره سطح وابستگی یک عبارت حمل می‌نمایند. به بیان دیگر، رده‌ها بخشی از ساختار نوعی سیستم هستند، اما لایه‌ها ابزاری ساختاری برای کنترل وابستگی‌ها محسوب می‌شوند.

در سیستم‌های نوع محض لایه‌ای، هر متغیر به‌صورت صریح به یک لایه نسبت داده می‌شود. این نشانه‌گذاری لایه‌ای، مشابه نشانه‌گذاری رده‌ای که پیش‌تر معرفی شد، صرفاً جنبهٔ تزئینی ندارد، بلکه در قواعد نوع‌دهی نقش اساسی ایفا می‌کند. قیود لایه‌ای تعیین می‌کنند که در چه شرایطی تشکیل انواع تابع وابسته، انتزاع و کاربرد مجاز است.

هدف از معرفی لایه‌ها، کاهش توان بیان سیستم نیست، بلکه فراهم‌سازی بستری است که در آن بتوان وابستگی‌های نوعی را به‌گونه‌ای کنترل‌شده بازآرایی کرد. این کنترل دقیق، امکان تعریف ترجمهٔ حذف وابستگی را فراهم می‌کند؛ ترجمه‌ای که نقش محوری در انتقال خواص نرمال‌سازی از سیستم‌های ساده‌تر به سیستم‌های نوع محض لایه‌ای ایفا می‌نماید.

\subsection{تعریف رسمی سیستم}

\begin{defn}[سیستم نوع محض لایه‌ای]
فرض کنید
$n$
یک عدد صحیح نامنفی باشد. یک سیستم نوع محض،
\lr{n}-لایه‌ای
است؛ هرگاه مشخصات آن به شکل زیر باشد
\begin{align*}
\mathcal{S} &= \{ s_i \mid i \in [n] \}\\
\mathcal{A} &= \{ (s_i,s_{i+1}) \mid i \in [n-1] \}\\
\mathcal{R} &\subset \{ (s,s',s') \mid (s,s') \in \mathcal{S} \times \mathcal{S} \}
\end{align*}
که در آن
\begin{equation*}
[n] = \{1,\dots,n\}
\end{equation*}
\end{defn}

\textbf{مثال}
\begin{enumerate}
\item تنها سیستم نوع محض ۰-لایه‌ای، سیستم نوع محض خالی است؛
\begin{equation*}
\mathcal{S}=\mathcal{A}=\mathcal{R}=\emptyset    
\end{equation*}
\item دو سیستم نوع محض ۱-لایه‌ای وجود دارد:
\begin{align*}
&\mathcal{S}=\{s_1\} \quad \mathcal{A}=\emptyset \quad \mathcal{R}=\emptyset\\
&\mathcal{S}=\{s_1\} \quad \mathcal{A}=\emptyset \quad \mathcal{R}=\{(s_1,s_1)\}
\end{align*}
\item سیستم‌های نوع محض ۲-لایه‌ای، معادل کل مکعب لامبدا هستند.
\end{enumerate}

\subsection{خواص}

در این بخش برخی خواص مهم فرانظری سیستم‌های نوع محض را بررسی می‌کنیم، و خواهیم دید که زیرکلاسی که در بخش قبل معرفی کردیم، چه ویژگی‌های کاربردی‌ای را در اختیار ما قرار می‌دهد.

\begin{defn}[سیستم نوع محض تابعی]
یک سیستم نوع محض  را تابعی گوییم، هرگاه
\begin{itemize}
\item اگر $(s,t) \in \mathcal{A}$ و $(s,t') \in \mathcal{A}$ آنگاه $t=t'$.
\item اگر $(s,t,u) \in \mathcal{R}$ و $(s,t,u') \in \mathcal{A}$ آنگاه $u=u'$.
\end{itemize}
\end{defn}

\begin{lmm}[یکتایی نوع]
در یک سیستم نوع محض تابعی، نوع هر عبارت نوع‌پذیر، یکتاست. یعنی
\begin{equation*}
\forall \Gamma, M, A, B.\quad (\Gamma \vdash M : A) \land (\Gamma \vdash M : B) \implies A =_\beta B
\end{equation*}
\end{lmm}

\begin{defn}[رده‌های بالا و پایین]
\hfill
\begin{itemize}
\item رده
$s$
یک رده بالا است هرگاه
\begin{equation*}
\nexists s' \in \mathcal{S}.\; (s,s') \in \mathcal{A}
\end{equation*}
\item رده
$s$
یک رده پایین است هرگاه
\begin{equation*}
\nexists s' \in \mathcal{S}.\; (s',s) \in \mathcal{A}
\end{equation*}
\end{itemize}
\end{defn}

\begin{lmm}[رده بالا]
برای هر زمینه
$\Gamma$،
متغیر
$x$،
عبارات
$A$ و $B$
و رده بالا
$s$،
گزاره‌های زیر برقرار هستند:
\begin{align*}
&\Gamma \nvdash s:A\\
&\Gamma \nvdash x:s\\
&\Gamma \nvdash AB:s
\end{align*}
\end{lmm}

\begin{defn}[ماندگاری]
یک سیستم نوع محض، ماندگار است؛ هرگاه هر سه گزاره زیر برای آن برقرار باشد:
\begin{itemize}
\item تابعی باشد.
\item اگر $(s,t) \in \mathcal{A}$ و $(s',t) \in \mathcal{A}$ آنگاه $s=s'$.
\item $\mathcal{R} \subset \{(s,s',s') \mid (s,s') \in \mathcal{S} \times \mathcal{S}\}$
\end{itemize}
\end{defn}

\begin{rem}
از اینجا به بعد، از نماد
$(s,s')$
به جای رابطه
$(s,s',s')$
استفاده می‌کنیم.
\end{rem}

\begin{defn}[اجتماع مجزا دو سیستم نوع محض]
اجتماع مجزا دو سیستم نوع محض
$\lambda\mathcal{S}$ و $\lambda\mathcal{S'}$،
که با نماد
$\lambda\mathcal{S} \sqcup \lambda\mathcal{S'}$
نشان داده می‌شود، یک سیستم نوع محض با مشخصات زیر است
\begin{align*}
\mathcal{S}_{\lambda\mathcal{S} \sqcup \lambda\mathcal{S'}} &= \mathcal{S}_{\lambda\mathcal{S}} \sqcup \mathcal{S}_{\lambda\mathcal{S'}}\\
\mathcal{A}_{\lambda\mathcal{S} \sqcup \lambda\mathcal{S'}} &= \mathcal{A}_{\lambda\mathcal{S}} \cup \mathcal{A}_{\lambda\mathcal{S'}}\\
\mathcal{R}_{\lambda\mathcal{S} \sqcup \lambda\mathcal{S'}} &= \mathcal{R}_{\lambda\mathcal{S}} \cup \mathcal{R}_{\lambda\mathcal{S'}}
\end{align*}
\end{defn}
